\chapter{Conclusioni}
\label{cap:conclusioni}

\section{Che cosa \`e stato realizzato}

Il progetto implementa il prodotto $Y \leftarrow AX$ per matrice densa e
multivettore stretto, distribuito su una griglia di processi bidimensionale con
MPI e accelerato all'interno del singolo processo con CUDA. Il codice \`e in C11
per la parte host e distribuita, in CUDA C++ per i kernel di device, e copre
tutti i requisiti della traccia elencati in \cref{tab:requisiti}.

Le scelte che hanno avuto l'effetto maggiore, e che sono anche quelle su cui
poggia l'analisi:

\begin{itemize}[nosep,left=0pt]
  \item \textbf{lo schema~A} ($i \to j \to c$), che legge ogni elemento di $A$
        una volta sola e lo riusa nei $k$ accumulatori in registro, portando
        l'intensit\`a aritmetica da $2/s$ a $2k/s$
        $\si{\flop\per\byte}$ --- il massimo ottenibile leggendo $A$ una volta;
  \item \textbf{la griglia 2D}, che riduce il volume comunicato da $\Theta(kM)$
        o $\Theta(kN)$ delle decomposizioni monodimensionali a
        $\Theta(k\sqrt{MN}/\sqrt{P})$;
  \item \textbf{il ciclo di vita a tre fasi} del kernel locale, che sposta fuori
        dal cronometro tutto ci\`o che riguarda $A$ --- inclusa la copia H2D, che
        altrimenti misurerebbe il PCIe invece della GPU --- rendendo la
        prescrizione della traccia una propriet\`a strutturale e non una
        sottrazione a posteriori;
  \item \textbf{una variabile per volta} fra le varianti di kernel, che \`e la
        condizione perch\'e ogni differenza misurata sia attribuibile.
\end{itemize}

\section{Che cosa dicono i dati finora}

\begin{daverificare}
Questa sezione va riscritta a raccolta ultimata. Al momento l'unica campagna
disponibile \`e C3, e il suo risultato principale \`e riportato qui sotto.
\end{daverificare}

Il nucleo di CPU raggiunge \SI{10.9}{\giga\flop\per\second} a $k = 8$ su
$\num{10000}\times\num{10000}$ in doppia precisione a rank singolo, e presenta
una discontinuit\`a netta fra $k = 8$ e $k = 20$: il costo per elemento di $A$
passa da \SI{1.47}{\nano\second} a \SI{7.68}{\nano\second}, mentre il lavoro
cresce solo di $2{,}5\times$. L'analisi di \cref{sec:c3} attribuisce il
fenomeno alla rilettura di $X$, che oltre una certa soglia esce dalla cache
privata del core, e propone tre controlli per falsificare l'ipotesi. La firma a
sostegno \`e un plateau di banda effettiva su $X$ intorno a
\SIrange{17}{21}{\giga\byte\per\second}, costante su cinque valori di $k$ che
coprono un fattore $3{,}2$.

Il modello roofline prevede inoltre che sulla GPU, in doppia precisione, il
nucleo sia compute-bound per ogni $k \ge 6$, con tetto a
\SI{350}{\giga\flop\per\second} indipendente da $k$, mentre in singola
precisione resti memory-bound per tutti i $k$ della traccia con tetto
proporzionale a $k$. \`E una previsione forte e falsificabile, e le campagne C1 e
C2 la mettono alla prova.

\section{Lavori futuri}
\label{sec:lavori-futuri}

\paragraph{Blocking del ciclo su $j$ sul nucleo di CPU.} \`E la conseguenza
diretta dell'analisi di \cref{sec:c3}: percorrere $X$ a fette che stiano nella
cache privata, riusando ogni fetta su tutte le righe di $A$ prima di passare alla
successiva. Non cambia lo schema, cambia solo l'ordine di visita, ed \`e lo stesso
principio che \kernelname{cuda\_warp\_smem} applica sulla GPU con la shared
memory. \`E implementato in \kernelname{omp\_scheme\_a\_tiled}, con il budget del
tile esposto come parametro di build perch\'e la soglia di cache \`e un'ipotesi da
misurare e non un dato del codice; misure preliminari fuori dal server
confermano che il segno del guadagno dipende da $n_{loc}\,k\,\mathrm{sizeof}$
confrontato con la cache disponibile, cio\`e che il tiling paga quando la fetta
locale di $X$ non ci sta e costa quando ci sta gi\`a.

\paragraph{Campagna di misura dei backend OpenMP.} La traccia ammette OpenMP
\emph{oppure} CUDA e il progetto implementa entrambe le alternative dietro la
stessa interfaccia \code{local\_gemm}: \kernelname{omp\_scheme\_a} parallelizza
il ciclo sulle righe di $Y$, che \`e l'unico dei tre cicli del nucleo a non
richiedere sincronizzazione, e \kernelname{omp\_scheme\_a\_tiled} vi aggiunge il
tiling di $X$ in cache descritto nel paragrafo precedente. Il codice \`e
validato (\code{make check-omp} ripete l'intera suite a 1, 2 e 4 thread, perch\'e
una corsa fra thread non si manifesta con un thread solo) e la campagna \`e
predisposta in \code{experiments/c7\_openmp.txt}; quel che resta \`e eseguirla sul
server e riportarne i numeri. Le tre righe da confrontare a parit\`a di core
occupati sono $20\times1$, $4\times5$ e $1\times20$ fra rank e thread: differiscono
solo per dove passa il confine fra i due livelli di parallelismo.

\paragraph{MPI CUDA-aware.} L'interfaccia \code{local\_gemm} \`e stata progettata
agnostica rispetto alla posizione della memoria proprio per consentirlo:
passando puntatori device direttamente a \code{MPI\_Bcast} e \code{MPI\_Reduce}
si eliminerebbero H2D di $X$ e D2H di $Y$ dal cammino di ogni invocazione, senza
riscrivere n\'e \code{mpi\_matmul} n\'e i kernel. Il primo passo \`e verificare il
supporto sul nodo di calcolo con \code{ompi\_info | grep cuda}.

\paragraph{MPS per la condivisione della GPU.} Con MPS attivo i contesti di pi\`u
rank non si alternerebbero sulla scheda (\cref{sec:serializzazione}), e i
numeri multi-rank su GPU tornerebbero interpretabili come tempo di kernel invece
che come throughput aggregato. \`E anche il modo per replicare, con analisi
esplicita della serializzazione, i benchmark multi-griglia su GPU --- un punto
che i lavori di confronto lasciano scoperto.

\paragraph{Profilazione con \code{ncu}.} Le ipotesi formulate nei capitoli
precedenti hanno tutte un contatore hardware che le decide: occupancy e banda L2
per \kernelname{cuda\_warp\_smem}, conflitti di banco per lo sweep su
\code{SMEM\_PAD}, registri per warp per \kernelname{cuda\_warp\_tiled}, kernel
interno selezionato per \kernelname{cublas}. \`E il passo che trasforma le
spiegazioni proposte in spiegazioni verificate.
